% ====================================================================
% JAISP paper -- target: ApJ. Itemized outline (v2, in SH style).
% ====================================================================
\documentclass[twocolumn]{aastex701}

\section{Introduction}
\begin{itemize}
    \item big surveys and the need of joint processing
    \item what is the state of AI and can it replace traditional
    \item growth interest on foundation models, while they are not an answer to everything.
    \item prior work positioning: AstroCLIP, Multimodal Universe, AstroPT, and classical forced photometry/astrometry -- where the novelty boundary is (pixel-level, multi-instrument, mixed-resolution fusion)
    \item start with already processed mosaics here, and emphasis on joint representation learning, but can be pushed earlier in the pipeline
    \item this paper
\end{itemize}

\section{Data}
\begin{itemize}
    \item multi-band multi survey imaging
    \item Rubin u g r i z y, DP1 ComCam coadds, patches, tracts, ...
    \item Euclid VIS and NISP Y J H -- imaging from DR1 MER mosaics, catalogues from Q1 MER (don't conflate Q1/DR1)
    \item MER vs NIR mosaics
    \item Gaia DR3 for PSF training and absolute astrometric reference
    \item Fields, area, depth, ...
    \item tiling and overlap: 50\% overlap, bit-identical overlap pixels, dedup needed, no leakage (patch-disjoint audit)
    \item variance maps vs MAD fallback
    \item figure 1: 10 band cutouts of a patch with RMS
\end{itemize}

\section{AI Architecture}
\begin{itemize}
    \item foundation with downstream heads
    \item Detection and Astrometry heads in this paper, in prep photometry and shape
    \item Foundation Model Evolution, where we started where we landed (latent/JEPA abandoned: supervised CNN 38 mas beat best JEPA 47 mas)
    \item foundation version insensitivity (v8$\approx$v9$\approx$v10): judge by OOD, not in-distribution probes
    \item Figure 2: Model architecture
    \item Validate learned foundation
    \item Figure 3: Contribution from different bands to the learned representation, resolution (linear probe, Euclid $R^2\sim0.8$)
\end{itemize}

\section{Detection}
\begin{itemize}
    \item architectures, centernet, stemcenternet (DETR-like was an abandoned foundation variant, not a current head)
    \item fine tuning with SEP, fine tuning with MER
    \item detection comparison on VIS (+ figure 4)
    \item injection test (+ figure 5)
    \item multi-band detection gain: VIS-only vs all-10-band completeness (+14 to +20 pp, 50\% depth $\sim$0.5 mag deeper)
    \item band-subset detection: NISP-only as high-z dropout proxy, Rubin-only $\sim$0\% (detector is Euclid-anchored)
    \item completeness and purity and how MER is not truth (+figure 6)
    \item what we gain and what we miss, and how to improve
\end{itemize}

\section{Astrometry}
\begin{itemize}
    \item Idea: use the detections + foundation to (a) measure positions in faint bands by transferring the sharp VIS localization, and (b) map the cross-survey concordance field -- not "beat the existing solution"
    \item How do Rubin and Euclid mosaics agree before any correction, given they are both GAIA aligned
    \item $\sim$40--50 mas raw offset is S/N dependent and dominated by per-source centering, not WCS distortion (+figure 7)
    \item including faint sources inflates the scatter but do they help with fitting the field (+figure 8)
    \item Head correcting for centering of faint things, and remeasuring the field (+figure 9)
    \item but the head has an intrinsic $\sim$17--18 mas faint-end ceiling; the classical joint 10-band fitter (6.7 mas at SNR 5) is the best position instrument
    \item injection truth test: real $\sim$2.5$\times$ gain at low SNR (VIS-localization transfer), bounded by the VIS floor
    \item centroid-emulation caveat: at the faint end the residual-vs-label metric measures emulation of the classical centroid, not truth -- why injection + Gaia are needed
    \item absolute validation against Gaia: frame tie $(+3.7,+1.9)$ mas, $4.4$--$4.8$ mas at $G<18.5$
    \item concordance field is a measurement (degree-scale $\sim$4--7 mas), not a correction: applying it does not improve held-out positions
    \item PINN vs HGP (agree at degree scale, sub-degree is solver noise)
    \item Ian's finding on centering dependence on a/b \todo{not in docs/notebooks -- needs a source/figure}
\end{itemize}

\section{OOD}
\begin{itemize}
    \item do results persist on out of distribution data (another field, EDF-S, no retraining)
    \item detection transfers nearly perfectly: completeness 89.6\% (OOD) vs 90.8\% (in-dist) -- figure 10
    \item astrometry: modest $\sim$5--15\% penalty (the scary $\sim$1.5$\times$ was a MAD-variance artifact, resolved with real variance) -- figure 11
    \item bright-star suppression reproduces (stronger, corr $-0.71$): origin is input normalization, not the field
    \item can be extended all sky
\end{itemize}

\section{Discussion}
\begin{itemize}
    \item Native rather than processed mosaics
    \item better PSF
    \item classical to inform model training, but faster inference once learned
    \item label noise limits in-distribution metrics; next head must beat 6.7 mas at SNR 5 (likely needs finer than 0.4"/px bottleneck)
    \item next paper
\end{itemize}
\end{document}
